\documentclass[11pt,a4paper]{article}
\usepackage[T1]{fontenc}
\usepackage[utf8]{inputenc}
\usepackage[main=italian,provide=*]{babel}
\usepackage{amsmath,amssymb}
\usepackage{geometry}
\geometry{margin=2.5cm}
\usepackage{booktabs}
\usepackage{longtable}
\usepackage{seqsplit}
\usepackage{array}
\usepackage{hyperref}
\hypersetup{colorlinks=true,linkcolor=blue,citecolor=blue,urlcolor=blue}

\title{Trimero ad anello (catena chiusa) — fasi e documenti prodotti}
\author{Note di lavoro — Tesi triennale, Samuele Schiavi\\ Relatore: Prof.\ Alessandro Chiesa}
\date{}

\begin{document}
\maketitle

\section{Vincolo di partenza}

Il trimero ad anello (triangolo isoscele) include tutte e tre le
interazioni $1\leftrightarrow2$ (legame di base, $J$),
$2\leftrightarrow3$ e $3\leftrightarrow1$ (legami laterali, $J'$),
a differenza della catena aperta dove il legame $1\leftrightarrow3$ è
assente per costruzione. La simmetria di scambio $1\leftrightarrow2$
(coppia di siti \emph{legata}) rende $S_{12}^2$ conservato e permette
la decomposizione di Kambe in forma chiusa, analogamente a come per la
catena la simmetria $P_{13}$ (coppia \emph{non} legata) conserva
$S_{13}^2$. Stesso percorso di lavoro seguito per il dimero e poi per
la catena: teoria esatta, stati espliciti, VQE — con l'aggiunta,
completata per l'anello, dell'analisi del termine DM.

\section{Stato di avanzamento}

\begin{longtable}{@{}p{1.6cm}p{2.0cm}p{7.3cm}p{2.2cm}@{}}
\toprule
Fase & Stato & Documenti prodotti & \\
\midrule
\endhead

1 & \textbf{Completata} &
\texttt{\seqsplit{teoria\_trimero\_isoscele.pdf}} -- teoria esatta:
identità $\sigma_1\cdot\sigma_2=2P_{12}-I$ e $S_{12}^2=I+P_{12}$
dimostrate da zero, condizione $[H,S_{12}^2]=0 \Leftrightarrow J'_{23}=J'_{31}$
provata esplicitamente (non assunta) via coniugazione per $P_{12}$,
decomposizione di Kambe con conteggio esplicito dei multipletti (metodo
del "peeling" invece della sola citazione di Clebsch--Gordan), spettro
chiuso a tre multipletti $A,B,C$, mappa dei segni completa nel piano
$(J,J')$ per quadrante, campo critico $b_c$ nei diversi regimi, prova
rigorosa che gli incroci sono veri (non anticrossing) via diagonalità
di $H$ nella base $|S_{12},S,M\rangle$, conseguenze pratiche per il VQE
(fondamentale quasi ovunque degenere, l'equilatero è il caso peggiore
per un test dell'ansatz).
&\\[4pt]
& &
\texttt{\seqsplit{animazione\_trimero\_isoscele.html}} -- visualizzazione
interattiva dello spettro in campo: geometria del triangolo, mappa di
fase 2D nel piano $(J,J')$, grafico dello spettro con fondamentale
colorato per blocco.
&\\[10pt]

2 & \textbf{Completata} &
\texttt{\seqsplit{derivazione\_stati\_kambe\_trimero.pdf}} -- derivazione
esplicita degli otto autostati nella base computazionale
$|q_1q_2q_3\rangle$, via operatore di abbassamento $S_-$ per il blocco
$A$ e ortogonalità per il blocco $B$ -- documento di riferimento poi
riusato, con lo stesso metodo, per derivare gli otto stati della
catena aperta.
&\\[10pt]

3 & \textbf{Completata} &
\texttt{\seqsplit{trimer\_ring\_exact.py}} -- benchmark esatto
(Hamiltoniana come \texttt{SparsePauliOp}, spettro analitico, stati di
Kambe verificati contro l'operatore Qiskit, proiettore sul
fondamentale, termini DM candidati), self-test a precisione macchina;
aggiornato in questa sessione per mediare $\langle M_z\rangle$ sul
sottospazio degenere invece che su un singolo autovettore arbitrario.
&\\[4pt]
& &
\texttt{\seqsplit{trimer\_ring\_exact.png}} -- spettro completo e
magnetizzazione del fondamentale, con distinzione esplicita fra il
punto convenzionale $b=0$ (fondamentale degenere) e l'incrocio fisico
vero a $b_c$ (stesso trattamento applicato poi alla catena).
&\\[4pt]
& &
\texttt{\seqsplit{verifica.py}} -- verifica numerica indipendente su
griglia: spettro completo (400 configurazioni casuali), energie dei tre
multipletti a $b=0$, campo critico (formula vs scansione numerica),
mappa dei segni (3721 punti), coordinate dei marker nelle figure --
script di controllo separato dal modulo principale, stesso ruolo di
\texttt{verifica\_catena\_preliminare.py} per la catena.
&\\[10pt]

4 & \textbf{Completata} &
\texttt{\seqsplit{vqe\_trimer\_ring.py}} -- primo modulo di produzione:
ansatz HA (hardware-efficient) e PMA "con branching" (blocco RBS su
tutti e tre i bond $12,23,31$, stato iniziale scelto fra i multipletti
di Kambe -- Tabella 2 di \texttt{derivazione\_stati\_kambe\_trimero.pdf}
-- a seconda che $b$ sia sopra o sotto $b_c(J,J')$); metodologia di
ottimizzazione a due stadi (multistart $R{=}6$, COBYLA + polish
L-BFGS-B) per mandato del relatore; self-test che verifica l'ansatz a
parametri nulli contro lo stato di Kambe atteso. \textbf{Rieseguito per
intero in questa sessione}: self-test a precisione macchina, script
principale completato senza errori, figura prodotta identica a quella
già in possesso.
&\\[4pt]
& &
\texttt{\seqsplit{vqe\_trimer\_ring\_spiegato.ipynb}} -- guida interattiva
al modulo precedente: principio variazionale su matrice $8\times8$,
blocco RBS, ansatz HA (9 par.) vs PMA (6 par., $K{=}1$), self-test
sull'ansatz a parametri nulli, VQE end-to-end su singolo punto e su
griglia ridotta ($\mathcal F=1$ ovunque, inclusa la degenerazione a
$b_c$), esplorazione interattiva del blocco scelto come stato iniziale
e del crollo del PMA vicino a $b_c$ con DM Opzione B ($D=0.15$) -- stessa
causa già isolata nel dimero (conservazione di $M$). Punto lasciato
esplicitamente aperto: un circuito nativo di preparazione per gli
stati-W (2 blocchi RBS in cascata, overlap $=1$ verificato a mano) è
stato derivato ma non ancora integrato in \texttt{vqe\_trimer\_ring.py}
al posto di \texttt{prepare\_state}. \textbf{Eseguito per intero in
questa sessione} (33 celle, zero errori): tutti i self-test a
precisione macchina, crollo di fidelity sotto DM a $b_c$
($\mathcal F=0.7262$) e recupero lontano da $b_c$
($\mathcal F=1.0000$) confermati.
&\\[4pt]
& &
\texttt{\seqsplit{vqe\_trimer\_ring\_W.py}} -- variante con il gate
$W_{ij}(\theta)$ di Crippa et al.\ al posto di RBS (mirror diretto
dell'ansatz "originale" del dimero, \texttt{vqe\_dimer.py}), stessa
struttura a branching; termine di paragone per isolare l'effetto della
sola scelta del gate $M$-conservante, coerente con
\texttt{scelta\_ansatz\_RBS\_vs\_W.md}. \textbf{Aggiunta in questa
sessione} una funzione di self-test dedicata (assente in origine),
speculare a quella di \texttt{vqe\_trimer\_ring.py} per uniformità fra i
due moduli -- verificata a precisione macchina ($2.22\times10^{-16}$).
&\\[4pt]
& &
\texttt{\seqsplit{vqe\_trimer\_ring\_W\_spiegato.ipynb}} -- guida
interattiva alla variante W: mirror diretto, sezione per sezione, di
\texttt{vqe\_dimer\_spiegato.ipynb}, qui su matrice $8\times8$; stessa
lettura del risultato centrale (HA sempre $\mathcal F=1$; PMA-W esatto a
$D=0$, crolla vicino a $b_c$ con DM Opzione B per la stessa ragione di
simmetria del dimero, $[H,S_{12}^2]\neq0$). \textbf{Eseguito per intero
in questa sessione} (29 celle, zero errori): stato iniziale verificato
(overlap $=1.000000$), PMA a 3 parametri esatto a $D=0$.
&\\[4pt]
& &
\texttt{\seqsplit{vqe\_trimer\_ring\_nobranch.py}} -- ansatz PMA
\textbf{senza} branching: stato iniziale fisso ($X$ su un solo qubit,
indipendente da $J,J',b$), la scelta del settore emerge
dall'ottimizzazione dei parametri anziché da una preparazione diversa
per regime; stesso conteggio di parametri della versione a branching a
$K{=}1$ (6), ma verificato numericamente \textbf{più robusto} al
termine DM (Opzione B, $D{=}0.15$, $b{=}b_c$: $\mathcal F=0.999898$ a
$K{=}1$, $\mathcal F=1$ esatto a $K{=}2$, contro il collasso di
fidelity della versione a branching vicino a $b_c$ -- vedi
\texttt{analisi\_dm\_trimero.pdf}). \textbf{Riprodotti in questa
sessione} tutti e tre i valori dichiarati nel docstring, coerenti anche
con \texttt{plateau\_dm\_ring\_cache.json}.
&\\[4pt]
& &
\texttt{\seqsplit{confronto\_ansatz\_entangler\_trimero\_anello.ipynb}} --
notebook \textbf{self-contained} finale (definizioni, calcolo e analisi
in un solo file, a differenza della fase esplorativa precedente divisa
su più moduli), rieseguito ed eseguito in questa sessione con cache
verificata: confronto sistematico di 10 ansatz VQE, fidelity con e
senza DM. Ansatz canonico per gate-count: \texttt{D: Ry-RBS-Ry} (9
parametri, 6 gate a due qubit, esatto anche con DM) -- non
\texttt{PMA-2q$\cdot$3} come indicato in una prima stesura di questo
stesso documento, correzione fattuale segnalata esplicitamente nel
notebook.
&\\[4pt]
& &
\texttt{\seqsplit{analisi\_espressivita\_PMA\_anello.ipynb}} -- indagine
mirata, stesso ruolo di \texttt{analisi\_rotazioni\_PMA.ipynb} (dimero)
ma su una domanda specifica dell'anello mai posta prima: la famiglia
"2q" (RBS su tutti e tre i legami, poi $R_y$ indipendenti) ha un
\textbf{tetto strutturale vero} $\mathcal F\approx0.99831$ vicino a
$b=0.05$ sotto DM Opzione B, identico per 6 e 9 parametri (verificato
con 60 restart e ottimizzazione diretta della fidelity, non
dell'energia); la famiglia "1q" -- inferiore a $D=0$ -- lo supera e
raggiunge $\mathcal F=1$ esatto a 9 parametri. Risultato
controintuitivo rispetto al caso $D=0$, con spiegazione strutturale
(le due famiglie esplorano sottospazi diversi dello spazio di Hilbert,
non annidati). Struttura e metodologia poi da replicare, se necessario,
per la catena aperta.
&\\[10pt]

5 & \textbf{Completata} &
\texttt{\seqsplit{analisi\_dm\_trimero.pdf}} -- analisi completa del
termine DM sul triangolo isoscele: fra due forme plausibili
(Opzione A: DM sui soli legami laterali; Opzione B: DM su tutti e tre i
legami, proporzionale a $J$), solo l'Opzione B apre effettivamente il
gap all'incrocio $b_c$ (dimostrato via la regola di non-incrocio di
von Neumann--Wigner: l'Opzione A preserva $S_{12}^2$ esattamente),
inclusa la ricostruzione dell'errore metodologico incontrato durante
l'analisi e di come è stato scoperto.
&\\

\bottomrule
\end{longtable}

\section{Ancora da fare}

Quantum simulation (Trotter) e correlazioni dinamiche sul trimero ad
anello: non ancora iniziate. Precedono cronologicamente lo stesso
lavoro sulla catena aperta (vedi
\texttt{\seqsplit{convenzioni\_qubit\_progetto.pdf}}, Parte 2, "Ancora da fare" e
\texttt{\seqsplit{fasi\_documenti\_catena\_aperta.pdf}}), dato che il relatore ha
chiesto di procedere prima con l'anello. La complicazione strutturale
attesa per Trotter sull'anello -- i tre legami di scambio condividono
qubit a coppie (i legami 12 e 23 condividono il qubit 2, e così via
ciclicamente), quindi $H_{\rm ex}$ non fattorizza esattamente e serve
una decomposizione di Trotter annidata -- è già registrata in
\texttt{log\_decisioni.md} come nota per quando questa fase verrà
affrontata.

\section{Documento di supporto trasversale}

\texttt{\seqsplit{convenzioni\_qubit\_progetto.pdf}} -- non specifico
dell'anello, ma aggiornato ad ogni fase del progetto: cataloga la
convenzione sito$\leftrightarrow$qubit di ciascun file (dimero, anello,
catena), organizzato nello stesso ordine cronologico della tesi, con il
"perché" di ogni scelta.

\section{Nota sulla numerazione delle fasi}

Come per la catena aperta, le fasi 1--5 qui sono state definite
\emph{a priori} come piano di lavoro (dimero $\to$ anello $\to$ catena
aperta), a differenza del dimero dove la suddivisione è stata
ricostruita \emph{a posteriori}. L'anello è l'unica delle tre strutture
per cui l'analisi del termine DM (Fase 5) risulta già completata: per
il dimero il DM è stato parte integrante fin dall'inizio (non una fase
a sé, vedi \texttt{fasi\_documenti\_dimero.pdf}), per la catena aperta
non è stato invece ancora richiesto.

\end{document}
